\documentclass[11pt,a4paper]{article}
\usepackage[utf8]{inputenc}
\usepackage[T1]{fontenc}
\usepackage{geometry}
\geometry{margin=2.5cm}
\usepackage{fancyhdr}
\usepackage{tcolorbox}
\usepackage{listings}
\usepackage{xcolor}
\usepackage{setspace}
\usepackage{titlesec}
\usepackage{enumitem}
\usepackage{eurosym}
\usepackage{hyperref}

% Entête et pied de page
\pagestyle{fancy}
\fancyhf{}
\rhead{Université de Lille - L2 Économie}
\lhead{TP Python — Corrigé détaillé}
\cfoot{\thepage}

% Hyperref (liens discrets)
\hypersetup{
  colorlinks=true,
  linkcolor=black,
  urlcolor=blue
}

% Style listings
\lstdefinelanguage{terminal}{
  morekeywords={ls,cd,pwd,mkdir,python,python3,code,dir,type,echo,py,conda,pip},
  sensitive=true
}
\lstset{
  basicstyle=\ttfamily\small,
  backgroundcolor=\color{gray!10},
  frame=single,
  language=Python,
  keywordstyle=\color{blue},
  commentstyle=\color{gray!70!black},
  stringstyle=\color{teal},
  showstringspaces=false,
  breaklines=true,
  columns=fullflexible
}

% Interligne
\onehalfspacing

% Couleur boîtes
\tcbset{
  colback=blue!5!white,
  colframe=blue!30!white,
  fonttitle=\bfseries,
  coltitle=black,
  boxrule=0.5pt,
  arc=2pt,
  left=6pt,
  right=6pt,
  top=4pt,
  bottom=4pt,
  before skip=10pt,
  after skip=10pt,
}

% Petites boîtes d'alerte/astuce
\newtcolorbox{AlerteBox}{colback=red!5!white,colframe=red!40!black,title=\textbf{Attention}}
\newtcolorbox{AstuceBox}{colback=green!5!white,colframe=green!40!black,title=\textbf{Astuce}}

\title{\textbf{Premiers pas en Python — Corrigé détaillé}}
\author{Université de Lille — L2 Économie}
\date{}

\begin{document}
\maketitle

\begin{tcolorbox}[title=\large Objectifs du TP (version corrigée pour révision)]
\begin{itemize}
  \item Vérifier votre installation de \textbf{Python 3} et de \textbf{VS Code}, savoir ouvrir un dossier et lancer un script.
  \item Comprendre la différence entre \textbf{terminal}, \textbf{interpréteur} et \textbf{script}.
  \item Manipuler les \textbf{types de base} (\texttt{int}, \texttt{float}, \texttt{str}, \texttt{bool}) et les \textbf{listes}.
  \item Maîtriser les \textbf{opérations arithmétiques}, l’\textbf{ordre des opérations} et les \textbf{conversions de types}.
  \item Savoir \textbf{indexer/slicer} des chaînes, \textbf{formater} l’affichage et raisonner avec des \textbf{booléens}.
\end{itemize}
\end{tcolorbox}

\section*{Installer Python et VS Code (rappel rapide)}
\begin{itemize}
  \item \textbf{Windows} : télécharger \href{https://www.python.org/downloads/}{python.org}, cocher \texttt{Add Python to PATH}, puis vérifier :
\begin{lstlisting}[language=terminal]
python --version
# ou
py -3 --version
\end{lstlisting}
  \item \textbf{macOS} : installer le \texttt{.pkg} puis :
\begin{lstlisting}[language=terminal]
python3 --version
\end{lstlisting}
  \item \textbf{Linux (Debian/Ubuntu)} :
\begin{lstlisting}[language=terminal]
sudo apt update
sudo apt install -y python3 python3-pip
python3 --version
\end{lstlisting}
  \item \textbf{VS Code} : installer \textit{Python (Microsoft)} et \textit{Jupyter} (facultatif).
\end{itemize}
\begin{AlerteBox}
Si \texttt{python3} ne fonctionne pas sous Windows, essayez \texttt{python}. Évitez Python~2 (\texttt{2.x}).
\end{AlerteBox}

\newpage
\begin{tcolorbox}[title=Interpréteur vs Script Python]
Quand vous travaillez avec Python, il existe deux manières de lancer du code :

\begin{enumerate}
  \item \textbf{L’interpréteur interactif} (dans le terminal, après avoir tapé \texttt{python3}).  
  Vous obtenez alors un prompt \texttt{>>>} qui vous permet de tester des instructions une par une.  
  Exemple :  
\begin{lstlisting}
>>> 2 + 3
5
>>> print("Bonjour")
Bonjour
\end{lstlisting}  
C’est idéal pour expérimenter rapidement et comprendre une erreur, comme une calculatrice évoluée.

  \item \textbf{Le script Python} (fichier \texttt{.py}).  
  Vous écrivez plusieurs lignes de code dans un fichier, par exemple \texttt{tp1.py}, puis vous l’exécutez en une fois :  
\begin{lstlisting}[language=terminal]
python3 tp1.py
\end{lstlisting}  
Le script conserve vos instructions : vous pouvez le modifier, le relancer, le partager. C’est ce mode qui doit être utilisé pour vos rendus.

\end{enumerate}

\textbf{Conseil pratique :} utilisez l’interpréteur pour tester et comprendre, puis recopiez vos solutions dans un fichier \texttt{.py} qui servira de rendu. Ainsi, vous profitez du meilleur des deux mondes.
\end{tcolorbox}

%==========================
\section*{Exercice 1 — Organisation (corrigé)}
\textbf{Commandes attendues :}
\begin{lstlisting}[language=terminal]
cd ~/Bureau           # (macOS/Linux) - sous Windows, aller dans le dossier choisi
mkdir Python_L2
cd Python_L2
mkdir TP1 data scripts notes
ls      # macOS/Linux
# dir   # Windows
\end{lstlisting}
Ouvrir le dossier dans VS Code :
\begin{lstlisting}[language=terminal]
code .
\end{lstlisting}
\textbf{Questions :}
\begin{itemize}
  \item Revenir en arrière : \texttt{cd ..}
  \item \texttt{ls} (Unix/macOS) $\approx$ \texttt{dir} (Windows) : liste le contenu.
\end{itemize}
\begin{AstuceBox}
\texttt{code .} ouvre le dossier courant dans VS Code (ne pas oublier l’espace avant le point).
\end{AstuceBox}

%==========================
\section*{Exercice 2 — Script vs Interpréteur (corrigé)}
\textbf{Script} \texttt{scripts/hello.py} :
\begin{lstlisting}
print("Bonjour Python !")
\end{lstlisting}
\textbf{Exécution dans le terminal} :
\begin{lstlisting}[language=terminal]
python3 scripts/hello.py
# Sortie attendue :
# Bonjour Python !
\end{lstlisting}
\textbf{Interpréteur (REPL)} :
\begin{lstlisting}[language=terminal]
python3
>>> print("Bonjour")
Bonjour
\end{lstlisting}
\textbf{À retenir} :
\begin{itemize}
  \item \textbf{Interpréteur} : tests rapides, ligne par ligne.
  \item \textbf{Script} : code réutilisable, partageable, versionnable.
\end{itemize}

%==========================
\section*{Exercice 3 — Python comme calculatrice (corrigé)}
\begin{lstlisting}
2 + 5      # 7
10 - 3     # 7
4 * 7      # 28
20 / 6     # 3.3333333333333335  -> division "réelle" (toujours float)
20 // 6    # 3                    -> division entière (quotient)
20 % 6     # 2                    -> reste
2 ** 5     # 32                   -> puissance
\end{lstlisting}
\textbf{Piège} : \texttt{/} renvoie \emph{toujours} un \texttt{float}, même si le résultat tombe « juste ».

%==========================
\section*{Exercice 4 — Ordre des opérations (corrigé)}
\begin{lstlisting}
2 + 3 * 4     # 14  (priorité à la multiplication)
(2 + 3) * 4   # 20  (parenthèses d'abord)
\end{lstlisting}
\textbf{Règle} : parenthèses $>$ puissances $>$ multiplications/divisions $>$ additions/soustractions.

%==========================
\section*{Exercice 5 — Chaînes de caractères (corrigé)}
\begin{lstlisting}
"Bonjour" + " le monde"   # "Bonjour le monde"
"Python" * 3              # "PythonPythonPython"
len("Université de Lille")# 21 (espaces comptent)
\end{lstlisting}
Indexation :
\begin{lstlisting}
s = "Python"
s[0]    # 'P'  (première lettre)
s[-1]   # 'n'  (dernière lettre)
\end{lstlisting}
\textbf{À retenir} : les chaînes se comportent comme des séquences de caractères.

%==========================
\section*{Exercice 6 — Variables (corrigé)}
\begin{lstlisting}
nom = "Alice"      # str
age = 22           # int
revenu = 1450.75   # float
etudiant = True    # bool
\end{lstlisting}
\textbf{Types} :
\begin{lstlisting}
type(nom), type(age), type(revenu), type(etudiant)
# <class 'str'>, <class 'int'>, <class 'float'>, <class 'bool'>
\end{lstlisting}
\textbf{int vs float} : \texttt{22} (entier), \texttt{22.0} (réel). \texttt{22 == 22.0} vaut \texttt{True} mais les \textbf{types} diffèrent.

%==========================
\section*{Exercice 7 — Conversions (corrigé)}
\begin{lstlisting}
age_str = "22"
int(age_str) + 3        # 25
str(22) + " ans"        # "22 ans"

int("3.14")             # ValueError (impossible de convertir "3.14" en int)
float("3.14")           # 3.14 (conversion OK en float)
\end{lstlisting}
\textbf{Leçon} : toujours convertir explicitement avant d’opérer.

%==========================
\section*{Exercice 8 — Affichage formaté (corrigé)}
\begin{lstlisting}
nom, age = "Alice", 22
print(nom, "a", age, "ans")     # ajoute des espaces automatiquement
print(f"{nom} a {age} ans")     # f-string : clair et puissant
revenu = 1450.756
print(f"Revenu: {revenu:.2f} €")# Revenu: 1450.76 €
\end{lstlisting}

%==========================
\section*{Exercice 9 — Listes (corrigé)}
\begin{lstlisting}
notes = [12, 15, 9]
notes[0]          # 12
notes.append(14)  # [12, 15, 9, 14]
notes[0] = 20     # [20, 15, 9, 14]
moyenne = sum(notes) / len(notes)  # 14.5
\end{lstlisting}
\textbf{Attention} : les listes sont \emph{mutables} (on peut modifier leurs éléments).

%==========================
\section*{Exercice 10 — Indexation de chaînes (corrigé)}
\begin{lstlisting}
mot = "Python"
mot[0]     # 'P'
mot[-1]    # 'n'
mot[0:3]   # 'Pyt'        (début inclus, fin exclue)
mot[2:]    # 'thon'       (du 3e caractère à la fin)
mot[-3:]   # 'hon'        (les 3 derniers caractères)
mot[::-1]  # 'nohtyP'     (chaîne inversée)
\end{lstlisting}

%==========================
\section*{Exercice 11 — Mise à jour de variable (corrigé)}
\begin{lstlisting}
x = 5
x = x + 1   # 6
x += 1      # 7 (raccourci équivalent)
\end{lstlisting}

%==========================
\section*{Exercice 12 — Booléens simples (corrigé)}
\begin{lstlisting}
x = 10
x > 5     # True
x == 10   # True  (égalité)
x != 8    # True  (différent)
x < 5     # False
\end{lstlisting}
\textbf{Important} : \texttt{=} affecte, \texttt{==} compare.

%==========================
\section*{Exercice 13 — Opérateurs logiques (corrigé)}
\begin{lstlisting}
a, b = True, False
a and b   # False
a or b    # True
not a     # False
\end{lstlisting}
\textbf{Tables de vérité} :
\begin{itemize}
  \item \texttt{and} : True si \emph{les deux} sont True.
  \item \texttt{or}  : True si \emph{au moins un} est True.
  \item \texttt{not} : inverse True/False.
\end{itemize}

%==========================
\section*{Exercice 14 — Booléens implicites (corrigé)}
\begin{lstlisting}
bool(0)        # False
bool(42)       # True
bool("")       # False   (chaîne vide)
bool(" ")      # True    (un espace = non vide)
bool([])       # False   (liste vide)
bool([1,2,3])  # True
\end{lstlisting}
\textbf{Règle} : les \emph{conteneurs vides} valent \texttt{False}, les non vides \texttt{True}.

%==========================
\section*{Exercice 15 — Mini-défis (corrigé)}
\begin{enumerate}
  \item Variables et opérations :
\begin{lstlisting}
x, y = 8, 3
print(x + y)   # 11
print(x - y)   # 5
print(x * y)   # 24
print(x / y)   # 2.6666666666666665
print(x // y)  # 2
print(x % y)   # 2
\end{lstlisting}

  \item \texttt{input()} :
\begin{lstlisting}
prenom = input("Votre prénom : ")
print("Bonjour", prenom)
# ou
print(f"Bonjour {prenom}")
\end{lstlisting}

  \item Listes :
\begin{lstlisting}
prenoms = ["Lucie", "Amine", "Sara"]
print(prenoms[1])  # "Amine"
len(prenoms)       # 3
\end{lstlisting}

  \item Types et conversions :
\begin{lstlisting}
"2" + 3          # TypeError
int("2") + 3     # 5
"2" * 3          # "222"
\end{lstlisting}
\end{enumerate}

%==========================
\section*{Messages d'erreur fréquents (et comment les lire)}
\begin{itemize}
  \item \texttt{SyntaxError: invalid syntax} — parenthèse, guillemet, deux-points manquants\dots
  \item \texttt{NameError: name 'x' is not defined} — variable non définie/typo/ordre d'exécution.
  \item \texttt{TypeError: can only concatenate str (not "int") to str} — mélange \texttt{str} et \texttt{int} sans conversion.
  \item \texttt{ModuleNotFoundError: No module named 'xxx'} — paquet non installé (\texttt{pip install xxx}) ou environnement incorrect.
\end{itemize}
\begin{AstuceBox}
Lisez le \textbf{dernier bloc} de l’erreur : il pointe l’endroit exact et la nature du problème.
\end{AstuceBox}

%==========================
\section*{Annexe — Lancer un script dans VS Code}
\begin{enumerate}
  \item Ouvrir un dossier (\texttt{Fichier} $\rightarrow$ \texttt{Ouvrir le dossier}).
  \item Créer \texttt{tp1.py} et y écrire du code.
  \item Installer l’extension \textit{Python (Microsoft)} si nécessaire.
  \item Sélectionner \textbf{l’interpréteur} Python (barre d’état en bas à droite).
  \item Lancer \texttt{Run} (triangle vert) ou utiliser le terminal intégré :
\begin{lstlisting}[language=terminal]
python3 tp1.py   # macOS/Linux
python tp1.py    # Windows
\end{lstlisting}
\end{enumerate}

\end{document}
